\uuid{mXDD}
\titre{Étude de fonction}
\auteur{bodin ; hanani ; mzari}
\organisation{exo7}
\licence{CC BY-NC-SA 4.0}
\datecreate{2020-01-20}
\niveau{L1}
\module{Analyse}
\chapitre{Fonctions circulaires et hyperboliques inverses}
\sousChapitre{Autre}
\difficulte{2}
\theme{Fonctions usuelles, étude de fonction}
\exo7qcmid{lille-1.fonctions-usuelles.016}

\contenu{
\begin{question}
\qclassification{126.00}
\qtags{motcle=étude de fonction}
Soit   $f(x)= \arctan x + \arctan (\frac{1}{x})$.  Quelles sont les assertions vraies ?

\begin{answers}
    \good{Le domaine de définition de $f$ est $\Rr^*$.}

    \bad{$f$ est une fonction constante.}

    \bad{$\forall x\in \Rr^*, \, f(x)=\frac{\pi}{2}$}

    \good{$\forall x>0, \, f(x)= \frac{\pi}{2}$ }

\end{answers}
\begin{explanations}
$f$ est définie sur $\Rr^*$,  dérivable sur  $\Rr^*$ et $f'(x)=0$, $\forall x \in \Rr^*$. Puisque $\Rr^* =]-\infty,0[\cup ]0,+\infty[$, $f$ est constante sur chaque intervalle. Or $f(1)=\frac{\pi}{2}$ et $f(-1)=-\frac{\pi}{2}$, on déduit que $f(x)=\left\{\begin{array}{cc}\frac{\pi}{2},& \mbox{si} \, \, x >0 \\ -\frac{\pi}{2},& \mbox{si} \,  x <0  \end{array}\right.$.
\end{explanations}
\end{question}
}
